Distributed publish--subscribe architectures---such as the Robot Operating System (ROS~2), Data Distribution Service (DDS), and MQTT---form the backbone of modern robotics, cyber-physical systems, and event-driven microservices. While asynchronous topic-based communication provides spatial and temporal decoupling, it obscures topological dependencies and failure propagation pathways, creating an \textbf{Architecture-Implementation Gap} where pristine component codebases can harbor severe system-level vulnerabilities. In this tool paper, we present \textbf{Software-as-a-Graph (SaG)}, an extensible, open-source model-driven engineering (MDE) framework for graph-based architectural modeling, diagnostic quality analysis, and counterfactually verified refactoring of publish--subscribe software systems.

SaG extracts typed, multi-layer multigraph models from Architecture-as-Code descriptors, folds source-level static code metrics into a \textbf{Code Quality Penalty (CQP)}, and computes an explainable, deterministic quality attribution model grounded in ISO/IEC 25010 and ISO/IEC 25019 standards across Reliability (Fault Tolerance and Availability) and Maintainability. Over these architectural models, SaG formalizes and evaluates a catalog of nineteen publish--subscribe anti-patterns using adaptive box-plot thresholds. To remediate diagnosed technical debt, the tool features a prescriptive engine that compiles findings into typed graph-mutation refactoring operators---logical topic splitting, physical host anti-affinity reallocation, and transport QoS contract hardening. Crucially, SaG enforces a closed-loop \textbf{``verify-before-transform''} discipline: candidate refactorings are counterfactually evaluated in an in-memory simulation sandbox, admitting only edits whose cascade-impact reduction exceeds stochastic seed noise across propagation thresholds ($\overline{\Delta I} > \kappa \sigma_{\text{seed}}$).

SaG is engineered as a decoupled hexagonal architecture featuring a dual-repository pattern: a graph database adapter (\texttt{Neo4jRepository}) for persistent querying and interactive exploration, and a high-performance in-memory adapter (\texttt{MemoryRepository}) that executes over 14,000 counterfactual simulation sweeps in under 3~hours. The tool provides multi-modal interfaces serving diverse stakeholder personas, including an interactive web dashboard and graph explorer, a Python SDK, a modular CLI suite, a REST API, and automated CI/CD regression gates. We evaluate SaG across eight multi-domain benchmark scenarios (29--520 vertices) and three transcribed real-world open-source systems (Autoware.Universe, Apollo Cyber RT, and an MQTT cloud mesh). We candidly report key empirical findings and calibration realities, including ranking saturation among structural predictors (NDCG@10 $\approx 0.82$--$0.86$, demonstrating that domain-grounded attribution is the primary value over scalar centrality), catalog discrimination improvements on real-world topologies ($\kappa = +0.296$ vs.\ chance on synthetic graphs), and a $71.0\%$ refactoring admission rate. Finally, we articulate five transferable lessons learnt spanning metamodel substrate adequacy, in-memory sandboxing, and simulation-in-the-loop verification to benefit the broader systems modeling community.
